% Section 16: Limitations — revised per supervisor review (2026-08-28)
% Reorganized from 7 subsections → 4 thematic paragraphs.
% E32/E33 CUAR inconsistency made prominent; explanation added here (§8 has the hook sentence).

\section{Limitations and Threats to Validity}
\label{sec:limitations}

\paragraph{Statistical power and sample sizes.}
Several layers are underpowered for the effect sizes they report. The live benchmark (E27) uses $n=100$ per system (CUAR upper CI 3.7\%); the escalation layer (E38) uses $n=30$ (upper CI 11.35\%); the clarification and normalization layers (E35, E37) use $n=45$ and $n=37$ with control arms that saturate, so those two layers support no directional claim and are excluded from the main narrative. E29's per-cell rates are exact multiples of 2.0\% despite a recorded $n=1{,}000$, implying an effective granularity near $n=50$; we report its rates without relying on its confidence intervals. All binary outcomes in the main battery use 95\% Wilson intervals, but for small-$n$ layers the intervals are wide and must be read as such. The 3.8~ms verification latency and the 546.0~ms and 1.2~ms values recur across layers with p50~$=$~p95 exactly; these are assumed constants in the results file, not measured percentiles.

\paragraph{Confounds in individual layers.}
E38 assigns a different model to each arm (GPT-4o-Mini / Llama-3.1-8B / Gemini-2.5-Flash-Lite), so the escalation headline is model-confounded; a same-model replication is required. E38 also records a post-hoc classifier fix with pre-fix numbers discarded. The E29 parametric layer's ``BAA'' row is invariant across modes by construction (the constraint overrides the mode), so it is one measurement, not four. E34's L4 and L5 differ only in latency/cost; the safety gain of the 11-slot verifier is established by E28's B5-vs-B6 contrast instead. Two exploratory layers (E32, oracle-vs-predicted routing; E33, high-confidence wrong metadata) recorded non-zero CUAR values for BAA (11.0\% and 10.0\%) that appear to contradict the 0.0\% main-battery headline. As noted in Section~\ref{sec:results_authority_safety}, these layers tested routing metadata perturbation---a scenario where CUAR is not identically applicable to evidence certification---and both are labelled exploratory with unresolved audit status rather than hidden.

\paragraph{Scope: what is not measured or claimed.}
We have no yield, income, or adoption data. Expert agronomist ratings (E13) measure advice quality, not field outcome. The network profiles (E14) are simulations, not field connectivity measurements; cost figures are scenario-based. The vision modality is currently whole-image classification-only; it does not perform localized bounding-box detection or spatial segmentation. Future work will integrate localized disease segmentation (SAM / YOLOv11-seg) to validate lesion coverage and severity before dosage calculation. English and other languages were not evaluated. The knowledge-graph traversal result (E19) is a proof of mechanism on a 23-node graph, not a scalability claim. The E24 maintenance result is a single-cluster case study.

\paragraph{Reproducibility and audit completeness.}
The experiment battery was produced in two phases with different audit depth. Layers E02--E26 carry a full verification block (self-checks, determinism re-run, trace check, real-application check, git commit, environment, seed). Layers E27--E39, which include the headline live-API benchmarks, do not yet carry that block in the registry; they are labelled \texttt{FROZEN\_AND\_VERIFIED} but are reproducibility-incomplete until their verification blocks are added. We state this openly rather than implying a uniform audit trail. Model and API drift, and regulatory change over time, remain ongoing validity threats that the architecture mitigates (versioned evidence, frozen thresholds) but cannot eliminate.